% Configurazione
\documentclass{article}

\usepackage{titling} % Required for inserting the subtitle
\usepackage{graphicx} % Required for inserting images
\usepackage{tabularx} % Per l'ambiente tabularx (tabelle)
\usepackage{calc} % Sempre per le tabelle
\usepackage[hidelinks]{hyperref} % Per i collegamenti ipertestuali, ad esempio sulla table of contents
\usepackage[italian]{babel} % Per la lingua italiana nelle scritte automatiche
\usepackage[table]{xcolor} % Per colorare il testo e le celle delle tabelle
\usepackage{colortbl} % Per colorare le celle delle tabelle
\usepackage{lipsum} % Per generare lorem ipsum
\usepackage[normalem]{ulem} % Per sottolineare il testo
\usepackage{array} % Per la visualizzazione fluttuante di array di domande e risposte
\usepackage{ragged2e} % Pacchetto necessario per \justifying che giustifica il testo di tabelle
\usepackage{tikz} % Per spostare elementi nel documento in modo facile e veloce
\usepackage[a4paper, top=2.5cm, bottom=2.5cm, left=2.5cm, right=2.5cm]{geometry} % Per i margini della pagina
\usepackage{fancyhdr} % Per l'intestazione e il piè di pagina
\usepackage{amsmath} % Per scrivere formule matematiche, in particolare per il pedice G

\newcommand{\ulhref}[2]{\href{#1}{\uline{#2}}} % Nuovo comando per sottolineare i link
\newcommand{\ulref}[1]{\uline{\ref{#1}}} % Nuovo comando per sottolineare i collegamenti a immagini e tabelle
\setlength{\parindent}{0pt} % Rimuove il rientro automatico dei paragrafi
\usetikzlibrary{calc} % Libreria per il calcolo delle coordinate di TikZ
\pagestyle{fancy} % Stile della pagina, per l'intestazione e il piè di pagina
\renewcommand{\footrulewidth}{0.4pt} % Inserimento della linea orizzontale in basso
\setlength{\headsep}{1.4cm} % Spazio tra l'intestazione e il testo

\graphicspath{ {immagini/} {../../shared/immagini/} }



% Struttura
\begin{document}

\thispagestyle{plain} % Niente intestazione e piè di pagina


\begin{tikzpicture}[remember picture, overlay]
    % Punto di partenza al centro orizzontale nella metà superiore
    \coordinate (top_center) at ($(current page.north)!0.3!(current page.south)$);

    % UniPD: Logo e descrizione
    \node at (top_center) [anchor=north, xshift=-3cm, yshift=4.85cm] 
        {\includegraphics[width=0.15\textwidth]{Logo Universita di Padova.png}};
    \node at (top_center) [anchor=north, xshift=1.7cm, yshift=4.5cm]
        {\textcolor{red}{\textbf{Università degli Studi di Padova}}};
    \node at (top_center) [anchor=north, xshift=1.7cm, yshift=4cm]
        {\textcolor{red}{Laurea: Informatica}};
    \node at (top_center) [anchor=north, xshift=1.7cm, yshift=3.5cm]
        {\textcolor{red}{Corso: Ingegneria del Software}};
    \node at (top_center) [anchor=north, xshift=1.7cm, yshift=3cm]
        {\textcolor{red}{Anno Accademico: 2024/2025}};

    % SWEg Labs: Logo e descrizione
    \node at (top_center) [anchor=north, xshift=-2.85cm, yshift=1.5cm] 
        {\includegraphics[width=0.16\textwidth]{Logo SWEg.png}};
    \node at (top_center) [anchor=north, xshift=1.7cm, yshift=0.5cm]
        {\textbf{Gruppo: SWEg Labs}};
    \node at (top_center) [anchor=north, xshift=1.7cm, yshift=0cm]
        {Email: \textsf{gruppo.sweg@gmail.com}};
\end{tikzpicture}


\vspace{10cm}

{
\centering
\Huge\bfseries Norme di progetto\par
\vspace{0.5cm}
\Large Versione 1.0.0\par
}

\vspace{2cm}
\begin{center}
\begin{tabular}{r|l}
    \textbf{Versione} & 1.0 \\
    \hline
    \textbf{Stato} & Approvato \\
    \hline
    \textbf{Uso} & Esterno \\
    \hline
    \textbf{Approvazione} & \emph{Tutto il gruppo} \\
    \hline
    \textbf{Redazione} & Federica Bolognini \\
    \hline
    \textbf{Verifica} & ....... \\
    \hline
    \textbf{Distribuzione} & \emph{SWEg Labs} \\
    \hline
    \textbf{Destinatari} & Prof. Tullio Vardanega \\
                         & Prof. Riccardo Cardin \\
\end{tabular}
\end{center}


\newpage
% Intestazione
\fancyhead[R]{\includegraphics[width=0.16\textwidth]{sweg_logo_sito_inverted.png}} % Immagine a destra

% Piè di pagina
\fancyfoot[C]{\thepage}                % Numero di pagina al centro

\pagenumbering{arabic} % Numerazione araba per il contenuto


Egregio prof. Vardanega, \\
Egregio prof. Cardin, \\



\end{document}
